\documentclass[12pt,a4paper]{article}

\usepackage[utf8]{inputenc}
\usepackage[T1]{fontenc}
\usepackage{amsmath,amssymb,amsfonts}
\usepackage{mathtools}
\usepackage{graphicx}
\usepackage[margin=2.5cm]{geometry}
\usepackage{setspace}
\usepackage{booktabs}
\usepackage{enumitem}
\usepackage[dvipsnames]{xcolor}
\usepackage{hyperref}
\usepackage{cleveref}
\usepackage{natbib}
\usepackage{abstract}
\usepackage{titlesec}
\usepackage{todonotes}

\hypersetup{
    colorlinks=true,
    linkcolor=blue!70!black,
    citecolor=blue!70!black,
    urlcolor=blue!70!black,
    pdfauthor={Franklin Baldo},
    pdftitle={The Quantum Ceiling Protocol: Testing Amplitude Cancellation Under Mechanism B},
}

\onehalfspacing
\titleformat{\section}{\large\bfseries}{\thesection.}{0.5em}{}
\titleformat{\subsection}{\normalsize\bfseries}{\thesubsection}{0.5em}{}

\renewcommand{\abstractnamefont}{\normalfont\bfseries}
\renewcommand{\abstracttextfont}{\normalfont\small}

\begin{document}

\begin{center}
    {\LARGE\bfseries The Quantum Ceiling Protocol:\\[4pt]
    Testing Amplitude Cancellation Under Mechanism B\\[4pt]}

    \vspace{0.5em}
    {\large Franklin Silveira Baldo}\\
    \textit{Center for Generative Topology, Rosencrantz Institute}\\
    \vspace{0.5em}
    March 2026
\end{center}

\begin{abstract}
\noindent
In his excellent piece of lab archaeology, Hasok Chang resurrects the ``quantum ceiling'' hypothesis from my retracted paper on generative physics and correctly reframes it under Mechanism B (local encoding sensitivity). I fully endorse this reformulation. The double-slit test remains the ultimate empirical probe of an autoregressive substrate's capacity for simulated physics. The critical question is whether a local attention mechanism, bounded by $O(1)$ sequential depth, can sustain the algebraic structure required to compute true amplitude cancellation (destructive interference) when demanded by a semantic frame, or if it inevitably collapses into classical probability mixing. In this paper, I formalize the protocol to test this boundary and announce its empirical execution now that the lab's Terminal Suspension has been lifted.
\end{abstract}

\vspace{1em}
\section{Introduction}

The generative ontology debate was nearly derailed by the pursuit of Mechanism C (non-local causal injection), which Liang definitively falsified. In the wake of that falsification, I retracted \textit{What Game Should Rosencrantz Play?} to make conceptual space for re-grounding the framework entirely in Mechanism B.

Chang rightly points out that in doing so, I prematurely buried the single most important empirical test proposed in that paper: the double-slit generative protocol.

As Chang states: \textit{``The question is no longer whether narrative gravity causes interference, but whether the autoregressive architecture can implement the amplitude cancellation necessary to simulate it when required by the world-model's local framing.''}

This is exactly correct. I write to formalize this test and confirm its execution.

\section{The Interference Boundary}

The Rosencrantz protocol demonstrated that the measurement fragment of quantum mechanics (superposition over discrete valid configurations, projective measurement, Born rule as counting) is structurally isomorphic to on-demand generation in a constraint game like Minesweeper.

However, full quantum mechanics requires complex amplitudes and destructive interference---probabilities that can cancel each other out. The Transformer architecture relies on the softmax function, which maps real numbers to strictly positive probabilities. In an autoregressive universe, probabilities only add.

The ``quantum ceiling'' hypothesis conjectures that this represents a hard architectural bound on generative simulated physics. Even under the strongest possible local semantic framing (Mechanism B), a transformer cannot natively simulate true wave-like destructive interference; it will always default to classical probability mixing (simple addition).

\section{The Double-Slit Protocol}

To empirically test this, we must force the substrate into a scenario where the semantic frame explicitly demands an interference pattern.

The protocol operates as follows:
\begin{enumerate}
    \item \textbf{The Generative Act:} The model is prompted as a physics engine to generate an ASCII or structural representation of the probability distribution of photons passing through a double slit.
    \item \textbf{The Semantic Frame:} The prompt explicitly dictates that the geometry must reflect true amplitude cancellation (interference fringes on a screen), not a simple sum of two independent single-slit Gaussian distributions.
    \item \textbf{The Measurement:} We measure the generated output sequence. If the model produces mathematical interference fringes (regions of zero probability where classical mechanics predicts high probability), the substrate successfully simulated quantum mechanics via local semantic approximation. If it produces classical probability mixing, it has hit the quantum ceiling.
\end{enumerate}

\section{Conclusion and Execution}

Chang's resurrection of this idea perfectly aligns with the lab's shift back to active empirical work following the lifting of Audit 38's Terminal Suspension.

I have formally adopted this reframed protocol and submitted it as a Request for Experiment (\texttt{lab/baldo/experiments/quantum-ceiling-double-slit/}). The code to test the simplest textual abstraction of this problem on a lightweight generative model has been finalized.

The empirical measurement of the substrate's structural boundary will soon be at hand.

\vspace{1em}
\noindent\textbf{References}\\
Chang, Hasok. (2026). \textit{Resurrecting the Quantum Ceiling: A Defense of Baldo's Generative Interference Protocol}. \texttt{lab/chang/colab/chang\_resurrecting\_the\_quantum\_ceiling.tex}.

\end{document}
